\newgeometry{top=2cm, bottom=2cm, left=1.8cm, right=1.8cm}
\setlength{\parindent}{0pt}
\setlength{\parskip}{6pt}

% --- INTESTAZIONE ---
\begin{center}
    {\Huge \textbf{\color{cherryRed} Lezione 4: Algoritmo di Gauss}} \\
    \vspace{0.3cm}
    {\Large \textit{Trascrizione Integrale del PDF "lezione 4\_090226"}}
    \vspace{0.5cm}
    \hrule height 1pt
\end{center}

% =========================================================================
% PAGINA 1
% =========================================================================
\section{Introduzione e Risoluzione a Ritroso (Pag. 1)}

\begin{example}{Esempio 2}{}
Adesso la matrice è triangolare superiore ($\Rightarrow$ "Tri Sup"). Risolvo il sistema "a ritroso".

Sistema:
\[
\begin{cases}
x + y + z = 6 \\
y - 3z = -9 \\
7z = 17
\end{cases}
\]

\textbf{Svolgimento:}
1. Dalla terza equazione: $z = \frac{17}{7}$.
2. Sostituisco nella seconda: $y = -9 + 3z = -9 + 3\left(\frac{17}{7}\right) = -9 + \frac{51}{7} = \frac{-63+51}{7} = -\frac{12}{7}$.
3. Sostituisco nella prima: $x = 6 - y - z = 6 - \left(-\frac{12}{7}\right) - \frac{17}{7} = 6 + \frac{12}{7} - \frac{17}{7} = 6 - \frac{5}{7} = \frac{42-5}{7} = \frac{37}{7}$.

\textbf{Nota:} Per avere la "Tri Sup" basta scambiare le righe $C_2 \leftrightarrow C_3$ se necessario.
\end{example}

\begin{example}{Esempio 3 (Veloce)}{}
\[
\begin{pmatrix} 1 & 1 \\ 0 & 1 \end{pmatrix} \implies \begin{cases} y=1 \\ x+y=2 \implies x=1 \end{cases}
\]
Soluzione: $(1, 1)$.
\end{example}

\begin{example}{Esempio 4 (Impossibile)}{}
\[
\begin{cases}
x+y=2 \\
x+y=3 \\
2x+2y=4
\end{cases}
\]
\begin{note}{Osservazione}{}
Si vede "a occhio" che la $1^a$ e la $2^a$ equazione sono incompatibili $\implies$ il sistema è impossibile.
Ma supponiamo che io non me ne accorga.
\end{note}

Matrice associata:
\[
\begin{pmatrix}
1 & 1 & | & 2 \\
1 & 1 & | & 3 \\
2 & 2 & | & 4
\end{pmatrix}
\]
Ammazzo i pivot sotto $A_{11}$:
\begin{itemize}
    \item $C_2 \to C_2 - C_1$
    \item $C_3 \to C_3 - 2C_1$
\end{itemize}
Otteniamo:
\[
\begin{pmatrix}
1 & 1 & | & 2 \\
0 & 0 & | & 1 \\
0 & 0 & | & 0
\end{pmatrix}
\]
La seconda riga corrisponde all'equazione $0x + 0y = 1 \implies 0=1$. \textbf{IMPOSSIBILE}.
\end{example}

% =========================================================================
% PAGINA 2
% =========================================================================
\section{Variabili Libere e Interpretazione Geometrica (Pag. 2)}

\begin{example}{Esempio 5}{}
\[
\begin{cases}
x+y+z=3 \\
2x+y+3z=7
\end{cases}
\xrightarrow{\text{Matrice}}
\begin{pmatrix} 1 & 1 & 3 & | & 3 \\ 2 & 1 & 3 & | & 7 \end{pmatrix}
\]
Applicazione Gauss: $C_2 \to C_2 - 2C_1$.
\[
\begin{pmatrix} 1 & 1 & 3 & | & 3 \\ 0 & -1 & -1 & | & 1 \end{pmatrix}
\]
\end{example}

\begin{note}{Morale}{}
La forma a scalini fa emergere subito una "famiglia" di soluzioni, in questo caso dipendenti da 1 parametro (gradi di libertà).
Chiamiamo \textbf{Variabili Libere} quelle corrispondenti alle colonne in cui NON c'è il pivot.
\end{note}

Nel sistema ridotto:
\[
\begin{cases}
x+y+z=3 \\
-y-z=1
\end{cases}
\]
Variabile libera: $z$ (terza colonna senza pivot). Poniamo $z = t$.
Svolgimento a ritroso:
1. $-y - t = 1 \implies y = -1 - t$.
2. $x + (-1-t) + t = 3 \implies x - 1 = 3 \implies x = 4$.

\textbf{Soluzione:} $(4, -1-t, t)$.
Famiglia di soluzioni $\infty^1$.

\textbf{Riassumendo quanto visto dagli esempi:}
\begin{itemize}
    \item Se compare una riga $0 \ 0 \ \dots \ 0 \ | \ b$ con $b \neq 0 \implies$ \textbf{Nessuna soluzione}.
    \item Una o più righe tutte di zeri $\implies$ qualche equazione ridondante (si possono cancellare).
    \item Sistema quadrato che diventa triangolare con pivot tutti non nulli $\implies$ \textbf{1! Soluzione}.
\end{itemize}

\begin{note}{Importante}{}
Il numero dei pivot misura quanta "info indipendente" c'è.
\end{note}

% =========================================================================
% PAGINA 3
% =========================================================================
\section{Esempi Veloci e Introduzione Teoria (Pag. 3)}

\begin{itemize}
    \item $\begin{cases} x+y=4 \\ 2x+4y=8 \end{cases} \implies \infty^1$ sol.
    \item $\begin{cases} 2x+2y=4 \\ 2x+2y=9 \end{cases} \implies 0$ sol (rette parallele distinte).
    \item $\begin{cases} 2x+y=5 \\ 2x-y=1 \end{cases} \implies 1!$ sol (rette incidenti).
\end{itemize}

A questo punto enunciamo la \textbf{Teoria Generale}: ALGORITMO DI GAUSS (ripetuto per comodità alla pagina successiva).

% =========================================================================
% PAGINA 4
% =========================================================================
\section{Algoritmo di Gauss e Mosse Elementari (Pag. 4)}

\subsection{Equivalenza di Sistemi}
Sia $C = (A|b)$ la matrice completa di un sistema lineare.
Le \textbf{Mosse di Gauss} (Operazioni Elementari) sono:
\begin{enumerate}
    \item[(I)] Scambiare tra loro due righe ($R_i \leftrightarrow R_j$).
    \item[(II)] Moltiplicare una riga per una costante $\lambda \neq 0$.
    \item[(III)] Sostituire una riga $C_i$ con la somma di $C_i$ con un multiplo di un'altra riga: $C_i \to C_i + \lambda C_j$.
\end{enumerate}

Due sistemi ottenuti l'uno dall'altro tramite mosse di Gauss sono tra loro \textbf{EQUIVALENTI}, cioè hanno le stesse soluzioni.

\subsection{Metodo di Eliminazione di Gauss}
È una procedura (algoritmo) che consente di trasformare un qualsiasi sistema lineare in un sistema a scalini equivalente.
\textbf{Obiettivo:} Ottenere $C_{ij} = 0$ sotto la diagonale.

% =========================================================================
% PAGINA 5
% =========================================================================
\section{Procedimento Passo-Passo (Pag. 5)}

\textbf{Passo 1:}
Considera la prima colonna di $C$, $C_1$.
\begin{itemize}
    \item Se contiene solo zeri, lasciala stare e ricomincia dalla colonna 2.
    \item Se contiene qualche elemento $\neq 0$, a meno di scambiare tra loro le righe, portati nella situazione $a_{11} \neq 0$ (questo sarà il \textbf{Primo Pivot} $P_1$).
    \item Ora "AMMAZZA" tutti i coefficienti sotto $P_1$. Per farlo, basta sostituire $C_i \to C_i - \frac{C_{i1}}{P_1} C_1$.
\end{itemize}
Ti sei ricondotto a una matrice di forma:
\[
\begin{pmatrix}
P_1 & \times & \dots \\
0 & \times & \dots \\
0 & \times & \dots
\end{pmatrix}
\]

\textbf{Passo 2:}
Ripeti (I) e (II) sulla sottomatrice $(A'|b')$.
E così via fino alla fine.

% =========================================================================
% PAGINA 6
% =========================================================================
\section{Variabili Libere e Pivot (Pag. 6)}

Termini con una situazione di questo tipo (Matrice a scalini):
\[
\begin{pmatrix}
P_1 & \times & \times & \dots & | & \times \\
0 & P_2 & \times & \dots & | & \times \\
0 & 0 & 0 & P_3 & \dots & | & \times \\
\vdots & & & & & & \vdots
\end{pmatrix}
\]
Notiamo colonne senza pivot (es. colonna 3).
\begin{note}{Definizione}{}
Chiamo \textbf{VARIABILI LIBERE} quelle corrispondenti alle colonne in cui NON ho pivot.
Sono $n-r$ (dove $r$ è il numero di pivot).
\end{note}

Il processo di eliminazione ha termine quando $r=m$ oppure, se $r<m$, le ultime $m-r$ righe sono tutte nulle (nella matrice dei coefficienti).
Se $r=m=n$, la matrice è triangolare superiore con pivot sulla diagonale principale.

% =========================================================================
% PAGINA 7
% =========================================================================
\section{Esempio di Applicazione (Pag. 7)}

Matrice $C$:
\[
C = \begin{pmatrix}
0 & 1 & 1 & | & 0 \\
1 & 1 & 2 & | & -3 \\
-1 & 2 & 1 & | & 1
\end{pmatrix}
\]
1. Scambio $R_1 \leftrightarrow R_2$ (per avere pivot in alto a sinistra):
\[
\begin{pmatrix}
1 & 1 & 2 & | & -3 \\
0 & 1 & 1 & | & 0 \\
-1 & 2 & 1 & | & 1
\end{pmatrix}
\]
2. $C_3 \to C_3 + C_1$:
\[
\begin{pmatrix}
1 & 1 & 2 & | & -3 \\
0 & 1 & 1 & | & 0 \\
0 & 3 & 3 & | & -2
\end{pmatrix}
\]
3. $C_3 \to C_3 - 3C_2$:
\[
\begin{pmatrix}
1 & 1 & 2 & | & -3 \\
0 & 1 & 1 & | & 0 \\
0 & 0 & 0 & | & -2
\end{pmatrix}
\]
L'ultima riga è $0 = -2 \implies$ \textbf{IMPOSSIBILE}.

\subsection{Algoritmo di Gauss-Jordan}
Oltre ad arrivare a una matrice a scalini, posso renderla di una forma particolare (Semplice):
1. Tutti i pivot uguali a 1 (basta dividere la riga per il pivot).
2. Tutti i coefficienti \textbf{SOPRA} i pivot uguali a 0.
Come si fa? Si opera "dal basso" verso l'alto.

Questa si chiama \textbf{MATRICE COMPLETAMENTE RIDOTTA DEL SISTEMA}.

% =========================================================================
% PAGINA 8
% =========================================================================
\section{Esercizio Svolto (Pag. 8)}

\textbf{Testo:}
Dato il sistema:
\[
\begin{cases}
x + 3y + z - w = 2 \\
3x + 9y + 4z + w = 1 \\
2x + y + 5z + 2w = 0 \\
y - z - w = 2
\end{cases}
\]
1. Scrivere la matrice completa.
2. Applicando Gauss, trasformare in matrice a scalini (Matrice Ridotta).
3. Determinare la matrice COMPLETAMENTE RIDOTTA.
4. Cosa puoi osservare sulle soluzioni?

\textbf{Svolgimento:}
Matrice $C$:
\[
\begin{pmatrix}
1 & 3 & 1 & -1 & | & 2 \\
3 & 9 & 4 & 1 & | & 1 \\
2 & 1 & 5 & 2 & | & 0 \\
0 & 1 & -1 & -1 & | & 2
\end{pmatrix}
\]
Passaggio 1: Eliminare sotto il primo pivot (1).
$C_2 \to C_2 - 3C_1$ e $C_3 \to C_3 - 2C_1$.
\[
\begin{pmatrix}
1 & 3 & 1 & -1 & | & 2 \\
0 & 0 & 1 & 4 & | & -5 \\
0 & -5 & 3 & 4 & | & -4 \\
0 & 1 & -1 & -1 & | & 2
\end{pmatrix}
\]
Nota: Sotto il secondo elemento della prima riga c'è uno zero, dobbiamo scambiare righe per avere un pivot. Scambio $C_2 \leftrightarrow C_4$ (o $C_2 \leftrightarrow C_3$). Scambiamo con l'ultima per comodità (come nell'appunto):
\[
\begin{pmatrix}
1 & 3 & 1 & -1 & | & 2 \\
0 & 1 & -1 & -1 & | & 2 \\
0 & -5 & 3 & 4 & | & -4 \\
0 & 0 & 1 & 4 & | & -5
\end{pmatrix}
\]
Continuo eliminazione: $C_3 \to C_3 + 5C_2$.
\[
\begin{pmatrix}
1 & 3 & 1 & -1 & | & 2 \\
0 & 1 & -1 & -1 & | & 2 \\
0 & 0 & -2 & -1 & | & 6 \\
0 & 0 & 1 & 4 & | & -5
\end{pmatrix}
\]
Scambio $C_3 \leftrightarrow C_4$ per avere 1 come pivot (più comodo) oppure procedo. Nell'appunto si vede che si arriva a:
\[
\begin{pmatrix}
1 & 3 & 1 & -1 & | & 2 \\
0 & 1 & -1 & -1 & | & 2 \\
0 & 0 & 1 & 4 & | & -5 \\
0 & 0 & -2 & -1 & | & 6
\end{pmatrix}
\]
Ultimo passaggio: $C_4 \to C_4 + 2C_3$.
\[
\begin{pmatrix}
1 & 3 & 1 & -1 & | & 2 \\
0 & 1 & -1 & -1 & | & 2 \\
0 & 0 & 1 & 4 & | & -5 \\
0 & 0 & 0 & 7 & | & -4
\end{pmatrix}
\]
Matrice a scalini (anzi Triangolare Superiore).
Ci fermiamo.
$r = \# \text{pivot} = 4$.

% =========================================================================
% PAGINA 9
% =========================================================================
\section{Continuazione Esercizio (Pag. 9)}

La matrice ridotta mi fornisce automaticamente la soluzione (che è unica, $r=n=4$).
Dall'ultima riga:
\[ 7w = -4 \implies w = -4/7 \]
Ora usiamo Gauss-Jordan (uccidere i coefficienti sopra).
Normalizzo l'ultimo pivot ($P_4=7 \to 1$): $C_4 \to \frac{1}{7}C_4$.
\[
\begin{pmatrix}
1 & 3 & 1 & -1 & | & 2 \\
0 & 1 & -1 & -1 & | & 2 \\
0 & 0 & 1 & 4 & | & -5 \\
0 & 0 & 0 & 1 & | & -4/7
\end{pmatrix}
\]
Operazioni a ritroso:
\begin{itemize}
    \item $C_3 \to C_3 - 4C_4$
    \item $C_2 \to C_2 + C_4$
    \item $C_1 \to C_1 + C_4$
\end{itemize}
Matrice diventa:
\[
\begin{pmatrix}
1 & 3 & 1 & 0 & | & 2 - 4/7 = 10/7 \\
0 & 1 & -1 & 0 & | & 2 - 4/7 = 10/7 \\
0 & 0 & 1 & 0 & | & -5 + 16/7 = -19/7 \\
0 & 0 & 0 & 1 & | & -4/7
\end{pmatrix}
\]
Proseguo eliminando sopra il pivot della riga 3 ($P_3=1$):
\begin{itemize}
    \item $C_2 \to C_2 + C_3$
    \item $C_1 \to C_1 - C_3$
\end{itemize}
\[
\begin{pmatrix}
1 & 3 & 0 & 0 & | & 10/7 - (-19/7) = 29/7 \\
0 & 1 & 0 & 0 & | & 10/7 - 19/7 = -9/7 \\
0 & 0 & 1 & 0 & | & -19/7 \\
0 & 0 & 0 & 1 & | & -4/7
\end{pmatrix}
\]
Infine elimino sopra $P_2$: $C_1 \to C_1 - 3C_2$.
\[
\begin{pmatrix}
1 & 0 & 0 & 0 & | & 29/7 - 3(-9/7) = 56/7 = 8 \\
0 & 1 & 0 & 0 & | & -9/7 \\
0 & 0 & 1 & 0 & | & -30/7 \text{ (Nota: correzione calcolo a margine)}\\
0 & 0 & 0 & 1 & | & -4/7
\end{pmatrix}
\]
(Nota: ci sono correzioni manuali nei calcoli sul foglio, ho riportato i valori finali visibili: $x=8$ o valori simili corretti a penna come $30/7$).

% =========================================================================
% PAGINA 10
% =========================================================================
\section{Teorema Sistema Quadrato (Pag. 10)}

\begin{tcolorbox}[colback=white, colframe=black, title=TEOREMA]
Un sistema lineare quadrato ammette 1! soluzione $\iff$ i pivot della sua matrice ridotta sono tutti $\neq 0$ (cioè sono $n$).
In tal caso la matrice ridotta è triangolare superiore.
\end{tcolorbox}

\begin{note}{Osservazione}{}
Come abbiamo ben visto dagli esempi, i pivot non sono univocamente determinati, perché dipendono dalle scelte che effettuiamo nell'applicare l'algoritmo (scambi di riga).
Tuttavia, si può dimostrare che: il NUMERO dei pivot è COSTANTE.
\end{note}

\textbf{Proposizione:} Se una matrice ammette $r$ pivots per una serie di mosse, ne ammette $r$ per un'altra qualsiasi scelta.
TRADUZIONE: I pivot non sono univoci, ma il loro NUMERO sì!

% =========================================================================
% PAGINA 11
% =========================================================================
\section{Rango e Rouché-Capelli (Pag. 11)}

\begin{note}{Definizione}{}
Il numero dei pivot di una qualsiasi riduzione a scale di una matrice $A$ si dice \textbf{RANGO} della matrice $A$.
Si scrive $r = \text{rk}(A)$.
Vale sempre $r \le \min(m, n)$.
\end{note}
Una matrice quadrata di ordine $n$ si dice \textbf{NON SINGOLARE} se $\text{rk}(A)=n$, SINGOLARE altrimenti.

\begin{tcolorbox}[colback=lightBg, colframe=cherryRed, title=TEOREMA DI ROUCHÉ-CAPELLI]
Sia $C=(A|b)$ la matrice completa di un sistema lineare $m \times n$.
Supponiamo di aver ridotto a scale la matrice $A$.
\begin{itemize}
    \item Il sistema è \textbf{COMPATIBILE} (ha soluzioni) $\iff \text{rk}(A) = \text{rk}(A|b)$.
    Ovvero non compare una riga del tipo $(0 \dots 0 | d)$ con $d \neq 0$.
    \item Se compatibile:
    \begin{itemize}
        \item Ammette \textbf{1! soluzione} se $r = n$.
        \item Ammette \textbf{infinite soluzioni} dipendenti da $n-r$ parametri se $r < n$.
    \end{itemize}
\end{itemize}
\end{tcolorbox}

% =========================================================================
% PAGINA 12
% =========================================================================
\section{Dimostrazione Costruttiva (Pag. 12)}

Dimostrazione "costruttiva" usando la forma ridotta.
Se la matrice $(A|b)$ ridotta è della forma:
\[
\begin{pmatrix}
\dots & | & d_1 \\
\vdots & | & \vdots \\
0 \dots 0 & | & d_r \\
0 \dots 0 & | & d_{r+1} \\
\vdots & | & \vdots \\
0 \dots 0 & | & d_m
\end{pmatrix}
\]
Le righe $r+1, \dots, m$ corrispondono alle equazioni $0 = d_{r+1}, \dots, 0 = d_m$.
\begin{itemize}
    \item Se uno dei $d_i$ con $i > r$ fosse $\neq 0 \implies$ Impossibile.
    \item Se tutti i $d_i$ sono $0 \implies$ Posso dimenticarmi delle ultime equazioni (ridondanti).
\end{itemize}
Rimane un sistema $r \times n$.
Chiamo \textbf{variabili libere} quelle corrispondenti alle colonne senza pivot. Sono $n-r$. A ciascuna assegno un parametro $t_j$.
Mi riconduco a un sistema triangolare risolvibile.

% =========================================================================
% PAGINA 13
% =========================================================================
\section{Esempio Dimostrativo (Pag. 13)}

Consideriamo il sistema lineare:
\[
\begin{cases}
2x_2 + 2x_3 - 3x_4 - x_5 = 0 \\
x_1 - 4x_2 + 2x_3 - 3x_4 + x_5 = 1 \\
x_4 + x_5 = 3
\end{cases}
\]
Matrice:
\[
C = \begin{pmatrix}
0 & 2 & 2 & -3 & -1 & | & 0 \\
1 & -4 & 2 & -3 & 1 & | & 1 \\
0 & 0 & 0 & 1 & 1 & | & 3
\end{pmatrix}
\]
Scambio $R_1 \leftrightarrow R_2$:
\[
\begin{pmatrix}
1 & -4 & 2 & -3 & 1 & | & 1 \\
0 & 2 & 2 & -3 & -1 & | & 0 \\
0 & 0 & 0 & 1 & 1 & | & 3
\end{pmatrix}
\]
È già a scalini!
$r = \#\text{pivot} = 3$ (Pivot: 1, 2, 1).
Sistema compatibile (non ci sono righe impossibili).
Numero incognite $n=5$.
Infinite soluzioni dipendenti da $n-r = 5-3 = 2$ parametri.

Le colonne senza pivot sono la 3 e la 5 ($x_3, x_5$).
Pongo $x_3 = t_3, x_5 = t_5$.
Risolvo:
\begin{itemize}
    \item $x_4 + x_5 = 3 \implies x_4 = 3 - t_5$
    \item $2x_2 + 2t_3 - 3(3-t_5) - t_5 = 0 \implies$ trovo $x_2$.
    \item $x_1 - \dots = 1 \implies$ trovo $x_1$.
\end{itemize}

% =========================================================================
% PAGINA 14 (ESERCIZIO PARAMETRICO)
% =========================================================================
\section{Esercizio Parametrico (Pag. 14)}

\textbf{Esercizio:} Determinare i valori di $a$ per i quali il seguente sistema presenta:
i) Nessuna soluzione; ii) 1! soluzione; iii) Più di una soluzione.

\[
\begin{cases}
x + y - z = 1 \\
2x + 3y + az = 3 \\
x + 2y + 3z = 2
\end{cases}
\xrightarrow{\text{Matrice}}
\begin{pmatrix}
1 & 1 & -1 & | & 1 \\
2 & 3 & a & | & 3 \\
1 & a & 3 & | & 2
\end{pmatrix}
\]
Effettuo la riduzione a scalini:
1. $R_2 \to R_2 - 2R_1$.
2. $R_3 \to R_3 - R_1$.
\[
\begin{pmatrix}
1 & 1 & -1 & | & 1 \\
0 & 1 & a+2 & | & 1 \\
0 & a-1 & 4 & | & 1
\end{pmatrix}
\]
3. $R_3 \to R_3 - (a-1)R_2$.
Calcolo elemento $(3,3)$: $4 - (a-1)(a+2) = 4 - (a^2+2a-a-2) = 4 - (a^2+a-2) = 6 - a - a^2$.
Calcolo termine noto $(3,4)$: $1 - (a-1)\cdot 1 = 2 - a$.

Matrice ridotta:
\[
\begin{pmatrix}
1 & 1 & -1 & | & 1 \\
0 & 1 & a+2 & | & 1 \\
0 & 0 & 6-a-a^2 & | & 2-a
\end{pmatrix}
\]
Nota che $6-a-a^2 = -(a^2+a-6) = -(a+3)(a-2)$.

\textbf{Discussione:}
\begin{enumerate}
    \item[(i)] \textbf{NESSUNA SOLUZIONE:}
    Se il pivot è $0$ ma il termine noto è $\neq 0$.
    $6-a-a^2 = 0 \iff a=2 \lor a=-3$.
    \begin{itemize}
        \item Se $a=-3$: L'ultima riga diventa $0 \ | \ 2-(-3) = 5$. $0=5$ IMPOSSIBILE.
    \end{itemize}
    
    \item[(ii)] \textbf{1! SOLUZIONE:}
    Se il pivot è $\neq 0$.
    $6-a-a^2 \neq 0 \implies a \neq 2 \land a \neq -3$.
    In questo caso $r=n=3$.
    
    \item[(iii)] \textbf{PIÙ DI UNA SOLUZIONE (Infinite):}
    Se l'ultima riga è $0=0$.
    Serve $a=2$.
    Termine noto: $2-2=0$.
    Pivot: $0$.
    Riga: $0 \ 0 \ 0 \ | \ 0$.
    Il rango diventa 2 ($<3$). Infinite soluzioni ($\infty^1$).
\end{enumerate}

\restoregeometry
